\section{Saadud andmed}

\begin{table}[H]
    \centering
    \begin{tabular}{|c|c|c|c|c|c|c|c|}
    \hline
        $t$ (s) & $t^\text{o}$ (°C) & $P$ (Pa) & $C_{\si{CO_2}}$ (ppm) & $AH$ (RH\%) & $a_x$ (m/s$^2$) & $a_y$ (m/s$^2$) & $a_y$ (m/s$^2$) \\ \hline
        0 & 25,88 & 99673,14 & 0 & 30,97 & 9,09 & 1,19 & -2,84 \\ \hline
        0 & 25,88 & 99675,52 & 0 & 30,97 & 9,12 & 1,05 & -3,35 \\ \hline
        0 & 25,87 & 99677,61 & 0 & 30,97 & 6,57 & 4,37 & -5,72 \\ \hline
        0 & 25,85 & 99677,76 & 0 & 30,97 & 6,27 & 4,67 & -6,44 \\ \hline
        1 & 25,03 & 99692,55 & 0 & 30,97 & 2,52 & 7,87 & -1,87 \\ \hline
        2 & 25,02 & 99691,23 & 0 & 30,97 & 0,3 & 0,54 & 0,29 \\ \hline
        3 & 24,91 & 99758,25 & 0 & 30,97 & -0,76 & 12,82 & 1,4 \\ \hline
        4 & 24,92 & 99860,72 & 0 & 30,97 & -1,16 & 12,74 & 3,57 \\ \hline
        5 & 24,92 & 99934,59 & 0 & 30,97 & -0,99 & 11,51 & 4,39 \\ \hline
        6 & 24,91 & 100002,56 & 0 & 30,97 & -2,4 & 19,6 & 4,24 \\ \hline
        7 & 24,88 & 100125,05 & 0 & 30,97 & -1,66 & 19,6 & 2,23 \\ \hline
        8 & 24,87 & 100199,05 & 0 & 30,97 & -0,42 & 13,36 & 1,88 \\ \hline
        9 & 24,82 & 100274,17 & 0 & 30,97 & -0,62 & 12,22 & 1,72 \\ \hline
        10 & 24,75 & 100354,12 & 0 & 30,97 & -1,79 & 11,15 & 1,53 \\ \hline
        11 & 24,6 & 100443,68 & 0 & 30,97 & -1,46 & 12,06 & 1,4 \\ \hline
        12 & 24,72 & 100526,24 & 0 & 30,97 & -1,75 & 11,58 & 2,04 \\ \hline
        13 & 24,72 & 100629 & 0 & 30,97 & -1,42 & 12,37 & 1,47 \\ \hline
        14 & 24,6 & 100721,22 & 0 & 30,97 & -1,37 & 12,34 & 1,53 \\ \hline
        15 & 24,67 & 100807,77 & 0 & 30,97 & -1,54 & 15,34 & 1,78 \\ \hline
        16 & 24,68 & 100894,87 & 0 & 30,97 & -1,24 & 12,66 & 1,84 \\ \hline
        17 & 24,68 & 100988,65 & 0 & 30,97 & -1,36 & 14,45 & 1,53 \\ \hline
        18 & 24,67 & 101080,45 & 0 & 30,97 & -1,43 & 9,08 & 1,45 \\ \hline
        19 & 24,63 & 101179,03 & 0 & 30,97 & -1,27 & 14,38 & 2,2 \\ \hline
        20 & 24,64 & 101266,55 & 0 & 30,97 & -1,38 & 15,29 & 1,87 \\ \hline
        21 & 24,63 & 101349,02 & 0 & 30,97 & -1,42 & 11,9 & 1,55 \\ \hline
        22 & 24,54 & 101435,18 & 0 & 30,97 & -1,54 & 12,94 & 1,61 \\ \hline
        23 & 24,54 & 101521,02 & 0 & 30,97 & -1,47 & 13,99 & 1,9 \\ \hline
        24 & 24,4 & 101620,73 & 0 & 30,97 & -1,18 & 12,64 & 1 \\ \hline
        25 & 24,26 & 101722,98 & 0 & 30,97 & -1,32 & 13,78 & 1,42 \\ \hline
    \end{tabular}
    \caption{Võistlusel saadud andmed}
    \label{tab:6}
\end{table}

$AH$ = absoluutne õhuniiskus

\section{Tulemuste analüüs}

\begin{itemize}
    \item CanSati langemisel suureneb õhurõhk, kuna ülevasse minnes läheme lähemalt kosmosesse, kus rõhk puudub.
    \item CanSati langemisel temperatuur väheneb, kuna soojem õhk tõuseb ülevale, samuti võiks mõjutada ka suurunenud CanSati ventilatsioon.
    \item 20eurone andur Aliexpressilt ei ole usaldusväärne ega saa mõõta õhuniiskuse ja CO2 sisalduse.
    \item Langevari vähendas laskumiskiirenduse umbes 8 m/s$^2$ võrra.
    \item Oodatud kiirus oli 8 m/s, tegelik kiirus tuli 6,8 m/s.
    \item CanSati kokkupõrkamisel maaga lendasid patareid välja, mis peatas andmete saatmist.

\end{itemize}

\section{Õppetunnid}

\begin{itemize}
    \item Saime aru, et meeskonna koostöö on päris raske.
    \item Saime teada, kuidas programmeerida, luua 3D mudeleid ja õmmelda langevarju.
    \item Õppisime, kuidas teha aruannet.
    \item Õppisime ületada elu raskusi.
    \item Saime aru, et ei tohi säästa raha andurite ostmisel.
\end{itemize}